\documentclass[../../main]{subfiles}

\begin{document}

\section{Relações de Equivalência}
\label{sec:matematica:relacoes:relacoes-de-equivalencia}

\begin{defn}[Relação de Equivalência]
  \label{def:matematica:relacoes:equivalencia}

  Seja (A) um conjunto e seja (R) uma relação binária sobre (A).

  Dizemos que (R) é uma relação de equivalência quando satisfaz:

  \begin{enumerate}
    \item reflexividade;

      ```
    \item simetria;

    \item transitividade.
      ```

  \end{enumerate}

  Isto é,

  [
    (\forall x \in A)
    (xRx),
  ]

  [
    (\forall x,y \in A)
    (xRy \Rightarrow yRx),
  ]

  e

  [
    (\forall x,y,z \in A)
    \bigl(
      (xRy \land yRz)
      \Rightarrow
      xRz
    \bigr).
  ]

\end{defn}

\begin{defn}[Classe de Equivalência]
  \label{def:matematica:relacoes:classe-equivalencia}

  Seja (R) uma relação de equivalência sobre (A).

  A classe de equivalência de um elemento (a \in A) é o conjunto

  [
    [a]
    =
    {
      x \in A
      \mid
      xRa
    }.
  ]

\end{defn}

\begin{prop}
  \label{prop:matematica:relacoes:classe-contem-elemento}

  Para todo (a \in A),

  [
    a \in [a].
  ]

\end{prop}

\begin{proof}

  Pela reflexividade,

  [
    aRa.
  ]

  Logo,

  [
    a \in [a].
  ]

\end{proof}

\begin{prop}
  \label{prop:matematica:relacoes:classes-iguais-ou-disjuntas}

  Se (R) é uma relação de equivalência sobre (A), então para
  quaisquer (a,b \in A),

  [
    [a]=[b]
  ]

  ou

  [
    [a]\cap[b]=\varnothing.
  ]

\end{prop}

\begin{proof}

  Suponha

  [
    [a]\cap[b]\neq\varnothing.
  ]

  Seja

  [
    x\in[a]\cap[b].
  ]

  Então

  [
    xRa
    \quad\text{e}\quad
    xRb.
  ]

  Pela simetria,

  [
    aRx.
  ]

  Pela transitividade,

  [
    aRb.
  ]

  Mostra-se então que

  [
    [a]\subseteq[b]
  ]

  e

  [
    [b]\subseteq[a].
  ]

  Portanto,

  [
    [a]=[b].
  ]

\end{proof}

\begin{prop}
  \label{prop:matematica:relacoes:particao}

  As classes de equivalência de uma relação de equivalência sobre (A)
  formam uma partição de (A).

\end{prop}

\begin{proof}

  Toda classe é não vazia pela reflexividade.

  Pela proposição anterior, duas classes são iguais ou disjuntas.

  Além disso,

  [
    A
    =
    \bigcup_{a\in A}[a].
  ]

  Logo as classes formam uma partição de (A).

\end{proof}

\begin{defn}[Conjunto Quociente]
  \label{def:matematica:relacoes:conjunto-quociente}

  Seja (R) uma relação de equivalência sobre (A).

  O conjunto das classes de equivalência é denominado conjunto
  quociente de (A) por (R) e é denotado por

  [
    A/R.
  ]

  Assim,

  [
    A/R
    =
    {
      [a]
      \mid
      a\in A
    }.
  ]

\end{defn}

\end{document}
